Read in Gauge Theory Gravity, the metrics below are derived metrics of synchronous gravitational fields on flat spacetime (in five dimensions by the Definition's $d$-dimensional clause), and the Ricci tensors are the components of the Ricci maps of their field strengths (Fact \claimref{synthesis.imported-gtg-dictionary}(iv)); the check computes in those components, as follows. Curvature is computed with $\Gamma^{a}{}_{bc} = \frac12 g^{ad}(\partial_{b}g_{dc} + \partial_{c}g_{db} - \partial_{d}g_{bc})$, $R^{a}{}_{bcd} = \partial_{c}\Gamma^{a}{}_{db} - \partial_{d}\Gamma^{a}{}_{cb} + \Gamma^{a}{}_{ce}\Gamma^{e}{}_{db} - \Gamma^{a}{}_{de}\Gamma^{e}{}_{cb}$, $\Ric_{bd} = R^{a}{}_{bad}$; this formula gives $\Ric = -\Lambda g$ on de Sitter in signature $(1,3)$, so the store's Ricci tensor is its negative, the calibration of Proposition \claimref{dark-energy.frozen-point-datum-kasner-wall}(ii), in which a massless scalar enters as $\Ric = -8\pi G\,d\phi\otimes d\phi$ (the focusing sign; Proposition \claimref{synthesis.unit-boost-residue-needs-nonpositive-exponent}(iii) uses the same normalisation, kinetic term $\frac12(\partial\phi)^{2}$).

\emph{(1) Reduction.} For the synchronous field $\hb^{-1}(e_{0}) = e_{0}$, $\hb^{-1}(e_{a}) = T^{P_{a}}e_{a}$ with $n$ spatial directions (derived metric $dT^{2} - \sum_{a=1}^{n}T^{2P_{a}}(dx^{a})^{2}$) the computation of Proposition \claimref{dark-energy.frozen-point-datum-kasner-wall}(i), (ii), which is plane by plane and holds for any $n$, gives the store-sign Ricci map $T^{2}\mathcal{R}(e_{0}) = (\sum_a P_a^{2} - \sum_a P_a)\,e_{0}$ and $T^{2}\mathcal{R}(e_{a}) = P_{a}(\sum_b P_b - 1)\,e_{a}$ (in components $T^{2}\Ric_{00}$ and $-T^{2}\Ric_{aa}$), so the field is Ricci-flat exactly when $\sum P_a = 1 = \sum P_a^{2}$; for $n = 4$ this is $P \in K_{5}$. The identity $G = T^{-P_4}g - T^{2P_4}(dx^4)^2$ is immediate from the definition of $g$. For $g = T^{P_4}(dT^{2} - \sum_{i}T^{2P_i}(dx^{i})^{2})$ the conformal factor $T^{P_4}$ and the exponents give, by the same diagonal computation, $\Ric_{ii} \propto (\sum_i P_i + P_4 - 1) = 0$ and, in the store's sign, $T^{2}\Ric_{TT} = -\frac32 P_4^{2}$ when $P \in K_{5}$, which equals $-8\pi G\,T^{2}(\partial_{T}\phi)^{2}$ for $\phi = \sqrt{3/16\pi G}\,P_4\log T$; the check computes all components exactly at three rational points of $K_{5}$ and shows that the relation fails at a point with $\sum P_a = 1$ off $K_{5}$. The wave equation: $\sqrt{|g|} = T^{(4P_4 + 2\sum_i P_i)/2} = T^{1 + P_4}$ and $g^{TT} = T^{-P_4}$, so $\sqrt{|g|}\,g^{TT}\partial_T\log T = 1$ and $\Box\phi = 0$ exactly. For the proper time, $d\tau = T^{P_4/2}dT$ gives $\tau = T^{k}/k$, $T = (k\tau)^{1/k}$, so $T^{P_4 + 2P_i} = (k\tau)^{2p_i}$ with $p_i = (P_i + P_4/2)/k$, and $\phi = \sqrt{3/16\pi G}\,(P_4/k)\log(k\tau)$, so $q = \sqrt{3/16\pi G}\,P_4/k$ and $8\pi Gq^{2} = \frac{3P_4^{2}}{2k^{2}}$. Then $\sum_i p_i = (1 - P_4 + \frac32 P_4)/k = 1$, and for any $P$ with $\sum_a P_a = 1$ and $S_2 := \sum_a P_a^{2}$,
\[
k^{2}\Bigl(\sum_i p_i^{2} + \tfrac32\tfrac{P_4^{2}}{k^{2}}\Bigr) = (S_2 - P_4^{2}) + P_4(1 - P_4) + \tfrac34P_4^{2} + \tfrac32P_4^{2} = S_2 - 1 + k^{2},
\]
so $\sum_i p_i^{2} + 8\pi Gq^{2} = 1 + (S_2 - 1)/k^{2}$, which is $1$ exactly on $K_{5}$.

\emph{(2) Bijection.} On $K_{5}$ every $|P_a| \le 1$, so $k \in [\frac12, \frac32]$ is positive and $\Pi$ is defined; by (1) it lands in $D_s$. On $D_s$, $|c| \le 1 < 2\sqrt{3/2}$, so $P_4 = c/(\sqrt{3/2} - c/2)$ is defined and is the unique solution of $c\,(1 + P_4/2) = \sqrt{3/2}\,P_4$; then $P_i = p_ik - P_4/2$ is forced by the formula for $p_i$, so $\Pi$ is injective. Given $(p, c) \in D_s$, define $P$ this way; then $\sum_a P_a = k\sum_i p_i - \frac32P_4 + P_4 = k - P_4/2 = 1$, and the identity of (1) with $\sum_i p_i^{2} + c^{2} = 1$ gives $S_2 = 1$, so $P \in K_{5}$ and $\Pi(P) = (p, c)$: $\Pi$ is surjective.

\emph{(3) The completed residue.} In $\Cl(1,4)$ the bivectors $B_a = e_0e_a$ ($a = 1, \dots, 4$) square to $-e_0^{2}e_a^{2} = 1$ and anticommute pairwise (for $a \neq b$, $B_aB_b = -e_0^{2}e_ae_b = -e_ae_b$ is antisymmetric in $a, b$), so $(\sum_a x_aB_a)^{2} = \sum_a x_a^{2}$. Hence $\hat B^{2} = \sum_i p_i^{2} + c^{2}$ and $B_5^{2} = \sum_a P_a^{2} = 1$; by Proposition \claimref{synthesis.unit-boost-residue-needs-nonpositive-exponent}(iii), given $\sum_i p_i = 1$, the Einstein-scalar equations at the wall hold exactly when $\sum_i p_i^{2} = 1 - 8\pi Gq^{2}$, that is when $\hat B^{2} = 1$, and then $1 - B_w^{2} = c^{2}$. Finally $k\,p_i = P_i + P_4/2$ is the displayed identity read coefficientwise on $B_1, B_2, B_3$; the term $\frac{P_4}{2}\sum_iB_i$ is the residue contributed by the factor $T^{P_4}$, since for $\Omega^{2}\tilde g$ with $\Omega^{2} = T^{P_4}$ the boost block of the Levi-Civita connection in the rescaled orthonormal frame acquires $\partial_T\log\Omega = P_4/(2T)$ on each spatial leg (the standard conformal change of the connection one-forms), and the factor $k$ is the conversion from $1/T$ in the frame of $\tilde g$ to $1/\tau$ in that of $g$ ($\tau^{-1}d\tau = k\,T^{-1}dT$). The check verifies the squares and the coefficient identity with GA.py in $\Cl(1,4)$ (with $e_4^{2} = -\frac32$ for $\hat B$, which absorbs $\sqrt{3/2}$ so that the coefficient $P_4/k$ is rational, and $e_4^{2} = -1$ for $B_5$). $\square$
